% Introduction paragraph for DDI Knowledge Graph project
% Covers: severity prediction, knowledge graph construction, DDI risk analysis, 
% alternative recommendations, and web application

\documentclass[11pt,a4paper]{article}
\usepackage[utf8]{inputenc}
\usepackage{geometry}
\geometry{margin=1in}

\title{\textbf{Introduction Paragraph for DDI Knowledge Graph Project}}
\author{}
\date{}

\begin{document}
\maketitle

\section*{Introduction}

While comprehensive drug-drug interaction (DDI) databases such as DrugBank catalog hundreds of thousands of interaction pairs, translating this information into actionable clinical guidance remains challenging: severity labels often lack evidence-based calibration, DDI data is siloed from mechanistic knowledge such as shared protein targets and metabolic pathways, and finding therapeutically equivalent alternatives with lower interaction burden requires considerable manual effort. Given the complexity of analyzing the entire DrugBank database, this work focuses on cardiovascular and antithrombotic drugs---a therapeutic domain where polypharmacy is common and DDI consequences are particularly severe. We present an integrated pipeline comprising: (1) a semantic severity recalibration methodology harmonizing DDI severity labels using natural language processing and clinical evidence hierarchies; (2) a multi-relational biomedical knowledge graph integrating drugs, proteins, pathways, diseases, side effects, and therapeutic categories from DrugBank, SIDER, and CTD; (3) graph-based polypharmacy risk assessment and ATC-guided alternative drug recommendations; and (4) a privacy-preserving web application with multi-modal drug input, automated severity classification, mechanistic explanations, and an integrated LLM chatbot---enabling clinicians to rapidly assess cardiovascular medication regimens and identify safer therapeutic alternatives.

\end{document}
